% figures/fig1/fig1.tex — auto-generated by src/adaptive_behavior/make_fig_tex.py
%
% Layout driven by exact PowerPoint coordinates from layout.tex.
% To regenerate this file after editing the pptx:
%   python -m adaptive_behavior.update_layout fig1
%   python -m adaptive_behavior.make_fig_tex fig1 --overwrite

% figures/fig1/layout.tex — auto-generated, do not edit by hand
% Regenerate: python -m adaptive_behavior.extract_pptx_layout figures/fig1/fig1.pptx --figure fig1

% Per-panel positions and sizes (exact PowerPoint measurements in mm)
\newlength{\FigIAX}\setlength{\FigIAX}{0.0mm}
\newlength{\FigIAY}\setlength{\FigIAY}{0.0mm}
\newlength{\FigIAW}\setlength{\FigIAW}{88.3mm}
\newlength{\FigIAH}\setlength{\FigIAH}{45.9mm}
\newlength{\FigIBX}\setlength{\FigIBX}{0.0mm}
\newlength{\FigIBY}\setlength{\FigIBY}{46.3mm}
\newlength{\FigIBW}\setlength{\FigIBW}{64.3mm}
\newlength{\FigIBH}\setlength{\FigIBH}{75.3mm}
\newlength{\FigICX}\setlength{\FigICX}{64.8mm}
\newlength{\FigICY}\setlength{\FigICY}{46.4mm}
\newlength{\FigICW}\setlength{\FigICW}{23.5mm}
\newlength{\FigICH}\setlength{\FigICH}{24.6mm}
\newlength{\FigIDX}\setlength{\FigIDX}{64.5mm}
\newlength{\FigIDY}\setlength{\FigIDY}{71.3mm}
\newlength{\FigIDW}\setlength{\FigIDW}{23.7mm}
\newlength{\FigIDH}\setlength{\FigIDH}{24.6mm}
\newlength{\FigIEX}\setlength{\FigIEX}{64.8mm}
\newlength{\FigIEY}\setlength{\FigIEY}{96.4mm}
\newlength{\FigIEW}\setlength{\FigIEW}{23.5mm}
\newlength{\FigIEH}\setlength{\FigIEH}{25.2mm}
\newlength{\FigIFX}\setlength{\FigIFX}{90.8mm}
\newlength{\FigIFY}\setlength{\FigIFY}{0.0mm}
\newlength{\FigIFW}\setlength{\FigIFW}{35.8mm}
\newlength{\FigIFH}\setlength{\FigIFH}{16.4mm}
\newlength{\FigIGX}\setlength{\FigIGX}{127.1mm}
\newlength{\FigIGY}\setlength{\FigIGY}{0.0mm}
\newlength{\FigIGW}\setlength{\FigIGW}{49.3mm}
\newlength{\FigIGH}\setlength{\FigIGH}{16.4mm}
\newlength{\FigIHX}\setlength{\FigIHX}{90.8mm}
\newlength{\FigIHY}\setlength{\FigIHY}{16.6mm}
\newlength{\FigIHW}\setlength{\FigIHW}{35.8mm}
\newlength{\FigIHH}\setlength{\FigIHH}{24.2mm}
\newlength{\FigIIX}\setlength{\FigIIX}{127.1mm}
\newlength{\FigIIY}\setlength{\FigIIY}{16.7mm}
\newlength{\FigIIW}\setlength{\FigIIW}{16.3mm}
\newlength{\FigIIH}\setlength{\FigIIH}{23.9mm}
\newlength{\FigIJX}\setlength{\FigIJX}{143.7mm}
\newlength{\FigIJY}\setlength{\FigIJY}{16.7mm}
\newlength{\FigIJW}\setlength{\FigIJW}{16.3mm}
\newlength{\FigIJH}\setlength{\FigIJH}{23.9mm}
\newlength{\FigIKX}\setlength{\FigIKX}{160.1mm}
\newlength{\FigIKY}\setlength{\FigIKY}{16.6mm}
\newlength{\FigIKW}\setlength{\FigIKW}{16.3mm}
\newlength{\FigIKH}\setlength{\FigIKH}{23.9mm}
\newlength{\FigILX}\setlength{\FigILX}{90.8mm}
\newlength{\FigILY}\setlength{\FigILY}{43.0mm}
\newlength{\FigILW}\setlength{\FigILW}{35.8mm}
\newlength{\FigILH}\setlength{\FigILH}{16.4mm}
\newlength{\FigIMX}\setlength{\FigIMX}{127.1mm}
\newlength{\FigIMY}\setlength{\FigIMY}{43.0mm}
\newlength{\FigIMW}\setlength{\FigIMW}{49.3mm}
\newlength{\FigIMH}\setlength{\FigIMH}{16.6mm}
\newlength{\FigINX}\setlength{\FigINX}{90.8mm}
\newlength{\FigINY}\setlength{\FigINY}{59.5mm}
\newlength{\FigINW}\setlength{\FigINW}{35.8mm}
\newlength{\FigINH}\setlength{\FigINH}{21.5mm}
\newlength{\FigIOX}\setlength{\FigIOX}{127.1mm}
\newlength{\FigIOY}\setlength{\FigIOY}{59.7mm}
\newlength{\FigIOW}\setlength{\FigIOW}{30.5mm}
\newlength{\FigIOH}\setlength{\FigIOH}{21.3mm}
\newlength{\FigIPX}\setlength{\FigIPX}{157.6mm}
\newlength{\FigIPY}\setlength{\FigIPY}{59.7mm}
\newlength{\FigIPW}\setlength{\FigIPW}{18.8mm}
\newlength{\FigIPH}\setlength{\FigIPH}{21.3mm}
\newlength{\FigIQX}\setlength{\FigIQX}{90.8mm}
\newlength{\FigIQY}\setlength{\FigIQY}{83.4mm}
\newlength{\FigIQW}\setlength{\FigIQW}{35.4mm}
\newlength{\FigIQH}\setlength{\FigIQH}{16.4mm}
\newlength{\FigIRX}\setlength{\FigIRX}{126.2mm}
\newlength{\FigIRY}\setlength{\FigIRY}{83.4mm}
\newlength{\FigIRW}\setlength{\FigIRW}{50.2mm}
\newlength{\FigIRH}\setlength{\FigIRH}{17.5mm}
\newlength{\FigISX}\setlength{\FigISX}{90.8mm}
\newlength{\FigISY}\setlength{\FigISY}{100.0mm}
\newlength{\FigISW}\setlength{\FigISW}{35.4mm}
\newlength{\FigISH}\setlength{\FigISH}{21.5mm}
\newlength{\FigITX}\setlength{\FigITX}{126.2mm}
\newlength{\FigITY}\setlength{\FigITY}{101.2mm}
\newlength{\FigITW}\setlength{\FigITW}{50.2mm}
\newlength{\FigITH}\setlength{\FigITH}{20.2mm}

% Derived group widths (bounding boxes of panel clusters)
\newlength{\FigIRightW}\setlength{\FigIRightW}{49.3mm}
\newlength{\FigILeftSubW}\setlength{\FigILeftSubW}{88.3mm}
\newlength{\FigIRightSubW}\setlength{\FigIRightSubW}{85.6mm}


\begin{figure}[H]
  \centering
  \setcounter{subfigure}{0}
  \resizebox{\linewidth}{!}{%
  \begin{tikzpicture}[inner sep=0pt, outer sep=0pt]

    \node[anchor=north west] at (\FigIAX, -\FigIAY) {%
      \begin{subfigure}[t]{\FigIAW}%
        \panel{A}{\panelgraphic[width=\linewidth]{figures/fig1/panels/fig1_A.pdf}}%
        \phantomsubcaption\label{fig:1-a}%
      \end{subfigure}};

    \node[anchor=north west] at (\FigIBX, -\FigIBY) {%
      \begin{subfigure}[t]{\FigIBW}%
        \panel{B}{\panelgraphic[width=\linewidth]{figures/fig1/panels/fig1_B.pdf}}%
        \phantomsubcaption\label{fig:1-b}%
      \end{subfigure}};

    \node[anchor=north west] at (\FigICX, -\FigICY) {%
      \begin{subfigure}[t]{\FigICW}%
        \panel{C}{\panelgraphic[width=\linewidth]{figures/fig1/panels/fig1_C.pdf}}%
        \phantomsubcaption\label{fig:1-c}%
      \end{subfigure}};

    \node[anchor=north west] at (\FigIDX, -\FigIDY) {%
      \begin{subfigure}[t]{\FigIDW}%
        \panel{D}{\panelgraphic[width=\linewidth]{figures/fig1/panels/fig1_D.pdf}}%
        \phantomsubcaption\label{fig:1-d}%
      \end{subfigure}};

    \node[anchor=north west] at (\FigIEX, -\FigIEY) {%
      \begin{subfigure}[t]{\FigIEW}%
        \panel{E}{\panelgraphic[width=\linewidth]{figures/fig1/panels/fig1_E.pdf}}%
        \phantomsubcaption\label{fig:1-e}%
      \end{subfigure}};

    \node[anchor=north west] at (\FigIFX, -\FigIFY) {%
      \begin{subfigure}[t]{\FigIFW}%
        \panel{F}{\panelgraphic[width=\linewidth]{figures/fig1/panels/fig1_F.pdf}}%
        \phantomsubcaption\label{fig:1-f}%
      \end{subfigure}};

    \node[anchor=north west] at (\FigIGX, -\FigIGY) {%
      \begin{subfigure}[t]{\FigIGW}%
        \panel{G}{\panelgraphic[width=\linewidth]{figures/fig1/panels/fig1_G.pdf}}%
        \phantomsubcaption\label{fig:1-g}%
      \end{subfigure}};

    \node[anchor=north west] at (\FigIHX, -\FigIHY) {%
      \begin{subfigure}[t]{\FigIHW}%
        \panel{H}{\panelgraphic[width=\linewidth]{figures/fig1/panels/fig1_H.pdf}}%
        \phantomsubcaption\label{fig:1-h}%
      \end{subfigure}};

    \node[anchor=north west] at (\FigIIX, -\FigIIY) {%
      \begin{subfigure}[t]{\FigIIW}%
        \panel{I}{\panelgraphic[width=\linewidth]{figures/fig1/panels/fig1_I.pdf}}%
        \phantomsubcaption\label{fig:1-i}%
      \end{subfigure}};

    \node[anchor=north west] at (\FigIJX, -\FigIJY) {%
      \begin{subfigure}[t]{\FigIJW}%
        \panel{J}{\panelgraphic[width=\linewidth]{figures/fig1/panels/fig1_J.pdf}}%
        \phantomsubcaption\label{fig:1-j}%
      \end{subfigure}};

    \node[anchor=north west] at (\FigIKX, -\FigIKY) {%
      \begin{subfigure}[t]{\FigIKW}%
        \panel{K}{\panelgraphic[width=\linewidth]{figures/fig1/panels/fig1_K.pdf}}%
        \phantomsubcaption\label{fig:1-k}%
      \end{subfigure}};

    \node[anchor=north west] at (\FigILX, -\FigILY) {%
      \begin{subfigure}[t]{\FigILW}%
        \panel{L}{\panelgraphic[width=\linewidth]{figures/fig1/panels/fig1_L.pdf}}%
        \phantomsubcaption\label{fig:1-l}%
      \end{subfigure}};

    \node[anchor=north west] at (\FigIMX, -\FigIMY) {%
      \begin{subfigure}[t]{\FigIMW}%
        \panel{M}{\panelgraphic[width=\linewidth]{figures/fig1/panels/fig1_M.pdf}}%
        \phantomsubcaption\label{fig:1-m}%
      \end{subfigure}};

    \node[anchor=north west] at (\FigINX, -\FigINY) {%
      \begin{subfigure}[t]{\FigINW}%
        \panel{N}{\panelgraphic[width=\linewidth]{figures/fig1/panels/fig1_N.pdf}}%
        \phantomsubcaption\label{fig:1-n}%
      \end{subfigure}};

    \node[anchor=north west] at (\FigIOX, -\FigIOY) {%
      \begin{subfigure}[t]{\FigIOW}%
        \panel{O}{\panelgraphic[width=\linewidth]{figures/fig1/panels/fig1_O.pdf}}%
        \phantomsubcaption\label{fig:1-o}%
      \end{subfigure}};

    \node[anchor=north west] at (\FigIPX, -\FigIPY) {%
      \begin{subfigure}[t]{\FigIPW}%
        \panel{P}{\panelgraphic[width=\linewidth]{figures/fig1/panels/fig1_P.pdf}}%
        \phantomsubcaption\label{fig:1-p}%
      \end{subfigure}};

    \node[anchor=north west] at (\FigIQX, -\FigIQY) {%
      \begin{subfigure}[t]{\FigIQW}%
        \panel{Q}{\panelgraphic[width=\linewidth]{figures/fig1/panels/fig1_Q.pdf}}%
        \phantomsubcaption\label{fig:1-q}%
      \end{subfigure}};

    \node[anchor=north west] at (\FigIRX, -\FigIRY) {%
      \begin{subfigure}[t]{\FigIRW}%
        \panel{R}{\panelgraphic[width=\linewidth]{figures/fig1/panels/fig1_R.pdf}}%
        \phantomsubcaption\label{fig:1-r}%
      \end{subfigure}};

    \node[anchor=north west] at (\FigISX, -\FigISY) {%
      \begin{subfigure}[t]{\FigISW}%
        \panel{S}{\panelgraphic[width=\linewidth]{figures/fig1/panels/fig1_S.pdf}}%
        \phantomsubcaption\label{fig:1-s}%
      \end{subfigure}};

    \node[anchor=north west] at (\FigITX, -\FigITY) {%
      \begin{subfigure}[t]{\FigITW}%
        \panel{T}{\panelgraphic[width=\linewidth]{figures/fig1/panels/fig1_T.pdf}}%
        \phantomsubcaption\label{fig:1-t}%
      \end{subfigure}};

  \end{tikzpicture}}% end resizebox

  % \caption{\textbf{Task design and behavior.}
  %   (\subref{fig:1-a}) schematic pdf
  %   (\subref{fig:1-b}) patch ethogram from example session
  %   (\subref{fig:1-c}) placeholder
  %   (\subref{fig:1-d}) p(stop) in first site
  %   (\subref{fig:1-e}) cumulative distribution of stops per patch
  %   (\subref{fig:1-f}) site level meme
  %   (\subref{fig:1-g}) example kinematics for site
  %   (\subref{fig:1-h}) example kinematics in 2d
  %   (\subref{fig:1-i}) intersite approach states
  %   (\subref{fig:1-j}) odorsite stopping states
  %   (\subref{fig:1-k}) odorsite leaving states
  %   (\subref{fig:1-l}) patch level meme
  %   (\subref{fig:1-m}) example kinematics for patch
  %   (\subref{fig:1-n}) example kinematics in 2d
  %   (\subref{fig:1-o}) average intersite traces for all mice
  %   (\subref{fig:1-p}) time to stop
  %   (\subref{fig:1-q}) session level meme
  %   (\subref{fig:1-r}) example kinematics in 3d
  %   (\subref{fig:1-s}) example states for runner
  %   (\subref{fig:1-t}) example states for stopper
  %   \emph{Insets}: Head-fixed mouse on wheel in three-screen VR environment with odor delivery; task graph showing mice running between patches, encountering odor sites, and choosing to stop (harvest) or run (leave); single-site velocity profile showing sustained stopping below velocity $v$ triggering a choice tone ($t$) followed by probabilistic reward ($d$); example velocity traces for short, long, and skipped patch visits.}
  \label{fig:1}
\end{figure}
